\begin{proposition}[Full-fibre cohomology and direct images]\hypertarget{tag-01F6}{}\label{cdp:st:C.3.11}
\StatementRecord{cdp:st:C.3.11}{proved}{01F6}{\SeriesRef{cdp:st:C.3.1}, \SeriesRef{cdp:st:C.3.12}, \SeriesRef{cdp:st:C.3.6}, \SeriesRef{cdp:st:D.2.7}, \SeriesRef{cdp:st:D.1.9}, \SeriesRef{cdp:st:C.3.7}, \SeriesRef{cdp:st:C.3.9}, \SeriesRef{cdp:st:C.3.10}, \SeriesRef{cdp:st:C.3.8}, \SeriesRef{cdp:st:D.1.12}}


Under the hypotheses of \ref{cdp:st:C.3.1}, a very general smooth fibre
is a two-torus or a primary or secondary Hopf surface of
algebraic dimension zero. The Inoue exclusion is \ref{cdp:st:C.3.12}.
Both Hopf
entries and the torus entry occur, by \ref{cdp:st:C.3.6} and \ref{cdp:st:D.2.7}.
For the entire smooth locus there are exactly the following
two cohomological alternatives:
\[
\begin{array}{c|c|c}
\text{smooth fibres}&
 (h^0(\OO),h^1(\OO),h^2(\OO))&
 \text{coherent direct images}\\ \hline
\text{two-tori}&(1,2,1)&
 R^2f_*\OO_Y\simeq\det R^1f_*\OO_Y\\
b_1=1,\ b_2=0&(1,1,0)&
 R^1f_*\OO_Y\simeq\OO_{\PP^1},\
 R^2f_*\OO_Y=0 .
\end{array}
\tag{C.PK.9}\hypertarget{tag-00VG}{}\label{cdp:eq:C.PK.9}
\]
The indicated cohomology dimensions hold on every
scheme-theoretic full fibre in the respective branch,
and both direct images satisfy base change there.
In the second branch
$H^q(Y,\OO_Y)=(\C,\C,0,0)$ and
$\Pic^0(Y)=\C\beta$ with the exponential convention
of \ref{cdp:st:D.1.9}.
\end{proposition}


\begin{proof}

\begin{enumerate}[leftmargin=*]
\item \ref{cdp:st:C.3.7}, \ref{cdp:st:C.3.9}, \ref{cdp:st:C.3.10} and \ref{cdp:st:C.3.12} remove the bielliptic,
Kodaira and Inoue types from \ref{cdp:st:C.3.1}; \ref{cdp:st:C.3.8} gives the Hopf algebraic
dimension. All smooth fibres are diffeomorphic over the
connected smooth base. In the torus branch they have
$b_1=4$, $b_2=6$, and are K\"ahler by the compact-surface
criterion. Their rational first Chern class and second
Chern number vanish by \ref{cdp:st:C.3.1}. The surface classification
then makes them two-tori: an effective nonzero canonical
divisor has strictly positive degree against a K\"ahler
form, whereas this rational first Chern class has zero
degree. The minimal K\"ahler surfaces of Kodaira dimension
zero with these Betti numbers are exactly the two-tori.
\ref{cdp:st:D.1.12} supplies this branch's full-fibre assertions.

\item In the Hopf branch every smooth fibre has
$b_1=1$ and $b_2=0$. A nonzero holomorphic two-form
is closed and has strictly positive integral against
its conjugate, so it defines a nonzero de Rham
cohomology class. Thus $b_2=0$ forces $p_g=0$.
Since $\chi(\OO)=0$ by \ref{cdp:st:C.3.1}, one has $q=1$.
By \cite[Theorem~2.2, pp.~4--5]{Viaclovsky2026Combs}, $\mathscr E=R^1f_*\OO_Y$ is a line bundle,
and its nowhere-zero transgression section $\beta$
from \ref{cdp:st:D.1.9} trivializes it. The coherent sheaf
$\mathscr G=R^2f_*\OO_Y$ has rank zero, hence finite
length. Hirzebruch--Riemann--Roch on $Y$ and coherent
Leray give the exact equality
\[
 0=\chi(Y,\OO_Y)
   =\chi(\OO_{\PP^1})-\chi(\mathscr E)+\chi(\mathscr G)
   =1-1+\operatorname{length}(\mathscr G).
\]
Therefore $\mathscr G=0$. Top-degree base change and
the parameter-multiplication exact sequence give
$h^2(\OO_{Y_c})=0$ and $h^1(\OO_{Y_c})=1$ on each
full fibre; \cite[Theorem~2.2, pp.~4--5]{Viaclovsky2026Combs} gives $h^0(\OO_{Y_c})=1$.
Leray now gives the stated global structure-sheaf
cohomology. Finally $H^1(Y;\Z)=0$ and the exponential
sequence identify $\Pic^0(Y)$ with
$H^1(Y,\OO_Y)=\C\beta$.

\end{enumerate}
\end{proof}